%% Shared manuscript body for Elsevier and arXiv formats.

\section{Introduction}
\label{sec:intro}

Financial fraud costs the global economy well over \$5 trillion per
year \citep{crowe2020}. On blockchain networks the problem is
particularly severe: Bitcoin transactions are pseudonymous, irreversible,
and publicly visible, making them a common vehicle for ransomware
payments, darknet commerce, and Ponzi schemes. Detection must happen
before confirmations accumulate and funds become unrecoverable.

Applying Graph Neural Networks (GNNs) to this problem has clear
theoretical motivation. Organised fraud does not operate through
isolated transactions. Money laundering requires layering funds across
multiple wallets; fraud rings share infrastructure; mixing services
leave topologically distinctive flow signatures. None of these signals
are visible when transactions are modelled as independent points, but
graph architectures that aggregate information across the transaction
network should recover them.

The Elliptic Bitcoin Dataset \citep{weber2019}, released in 2019,
labels 46,564 Bitcoin transactions across 49 temporal steps and has
become the standard benchmark for this hypothesis. Published results
have consistently supported the GNN advantage: GCN (F1\,=\,0.70,
\citep{weber2019}), augmented GCN (F1\,=\,0.74, \citep{alarab2020}),
GraphSAGE with self-supervised pre-training (F1\,=\,0.75,
\citep{lo2023}), and EvolveGCN (F1\,=\,0.77,
\citep{pareja2020evolvegcn}) all outperform feature-only baselines
reported alongside them.

This paper challenges that consensus by tightening the evaluation
protocol in two ways prior work did not. First, we enforce
\emph{inductive} evaluation: test-period nodes and their edges are
withheld from the graph during training, so the model must generalise
to unseen structure rather than propagating information through
test-time neighbourhoods. Every prior Elliptic study uses transductive
evaluation, which leaks ground-truth supervision forward through edges
that would not exist in a deployed system. Second, we report
performance on all 15 individual test steps rather than aggregating
over the full test window. That aggregation conceals a dramatic failure
mode: models are viable during steps 35--42 but nearly useless during
steps 43--49.

Under these stricter conditions, the results diverge sharply from the
published picture. A plain three-layer MLP achieves
F1\,=\,0.744\,$\pm$\,0.005 and outperforms every tested GNN.
GraphSAGE, the strongest graph model, reaches only
F1\,=\,0.697\,$\pm$\,0.009. For context, a Random Forest baseline
under the same temporal split achieves F1\,=\,0.674, which already
exceeds prior published Random Forest results (F1\,=\,0.64,
\citep{weber2019}) despite our stricter inductive protocol. Feature
importance analysis shows that 82.4\,\% of discriminative power comes
from local transaction features, the part that neither depends on nor
benefits from learned graph topology. Graph structure provides a
modest advantage during the first eight test steps when fraud patterns
resemble training, but that advantage reverses once the fraud
distribution shifts. Temperature scaling and threshold optimisation leave F1
unchanged, ruling out calibration as a remedy and pointing toward
periodic retraining as the only viable fix.

These findings are not an argument against GNNs in general. They are
a demonstration that benchmark conditions matter: the assumption of
distributional stability between training and test periods must be
verified explicitly, and inductive evaluation is the appropriate
default for any system modelling evolving transaction graphs.

This paper makes the following contributions:

\begin{enumerate}[leftmargin=*,label=\textbf{RQ\arabic*}]
  \item Which architecture--strategy combination performs best under
    inductive evaluation with statistical validation?
  \item How does performance vary across individual test time steps,
    and what drives the degradation in later steps?
  \item Do graph architectures provide discriminative benefit beyond
    node features alone, and if so, under what conditions?
  \item Can post-hoc calibration, temperature scaling or threshold
    tuning, recover performance lost to temporal shift?
  \item Does replacing the transaction graph with a constructed
    behavioural graph improve GNN performance, and by how much?
\end{enumerate}

\begin{itemize}[leftmargin=*]
  \item A statistically validated benchmark across 4 architectures
    $\times$ 3 imbalance strategies $\times$ 5 random seeds under
    inductive mini-batch evaluation, establishing the single-number
    performance ceiling for GNNs on this dataset.

  \item A systematic comparison of five graph construction strategies
    (original transaction graph, cosine similarity, feature k-NN,
    cross-timestep temporal, and augmented). The temporal graph
    construction improves GNN F1 by 89\,\% relative over the original
    transaction graph --- the most substantial intervention tested.

  \item A direct comparison against Random Forest and XGBoost under the
    same temporal split. Classical methods outperform single-run GNNs;
    the gap narrows under multi-seed averaging but does not close.

  \item An ablation decomposing GNN performance into contributions from
    graph structure and node content, showing that the original
    transaction graph is marginally harmful and that node content
    features carry virtually all discriminative power.

  \item A hybrid model evaluation (GNN embeddings $\to$ Random Forest)
    showing a 0.262 F1-point decline relative to RF on raw features,
    ruling out GNN representations as a feature-engineering shortcut.

  \item Feature importance analysis across RF (Gini), XGBoost (gain),
    and GNN (input-gradient saliency), showing convergence on the same
    feature subspace across three different model families.

  \item The first systematic per-timestep analysis on Elliptic, showing
    that aggregate test-set F1 conceals near-complete failure in steps
    43--49, with an impossibility result for post-hoc recalibration under
    those conditions.

  \item Code, checkpoints, and experiment scripts:
    \url{https://github.com/Saket-Maganti/gnn-fraud}.
\end{itemize}

\section{Related Work}
\label{sec:related}

\subsection{GNN-Based Fraud Detection}

The application of graph learning to transaction networks predates deep
learning. Early work used belief propagation and Markov random fields
to catch fraudulent accounts in payment graphs \citep{akoglu2015graph}.
Differentiable message-passing architectures made it possible to learn
the aggregation function end-to-end, improving substantially over those
earlier hand-crafted propagation rules.

On the Elliptic Dataset specifically, Weber et al.\ \citep{weber2019} provided the
founding GCN baseline (F1\,=\,0.70), noting that the 71 manually
engineered aggregate features in the dataset already provide a surrogate
for one round of message passing, so the additional gain from GCN over
a well-equipped MLP is modest. Pareja et al.\ \citep{pareja2020evolvegcn} introduced
EvolveGCN, which updates GCN weight matrices via a GRU across time
snapshots rather than fixing them once during training; on Elliptic
this yielded F1\,=\,0.77 but only with full GPU-based training over
all 34 training snapshots. Alarab et al.\ \citep{alarab2020} combined GCN with
feature engineering and class-weighted loss to achieve F1\,=\,0.74.
Lo et al.\ \citep{lo2023} applied GraphSAGE with structural graph augmentation
and self-supervised pre-training, reaching F1\,=\,0.75.

None of these works include an inductive evaluation protocol or report
per-timestep performance. The gap between what they measure and what
a deployed model would experience is the central motivation for our study.

Before GNN-based approaches, classical machine learning methods were the
standard for transaction fraud detection. Weber et al.\ \citep{weber2019}
reported Logistic Regression (F1\,=\,0.53) and Random Forest
(F1\,=\,0.64) as baselines on Elliptic under transductive evaluation.
Gradient-boosted trees such as XGBoost \citep{chen2016xgboost} and
LightGBM perform in a similar range on comparable imbalanced
classification tasks. These classical results matter for context:
as we show, our MLP under stricter evaluation outperforms published
GNN results that were compared against these weaker baselines.

Outside the Elliptic Dataset, GNN fraud detection has been applied to
e-commerce at scale \citep{wang2019semi}, financial services networks
\citep{liu2018geniepath}, and insurance claims \citep{pourhabibi2020fraud}.
A shared assumption in this work is distributional stability between
training and deployment, which holds in some settings but fails in ours.

\subsection{Handling Class Imbalance}

Fraud datasets are severely imbalanced: legitimate transactions
comfortably outnumber fraudulent ones by a factor of 7--10$\times$ in
Elliptic and often far more in production systems. Standard
cross-entropy loss allocates most gradient to the majority class,
leaving the model with little incentive to detect rare fraud.

Chawla et al.\ \citep{chawla2002smote} proposed SMOTE, which synthesises minority
samples by linearly interpolating between existing minority examples in
feature space. Zhao et al.\ \citep{zhao2021graphsmote} adapted this to graph data
by interpolating in embedding space and rewiring edges for synthetic
nodes. Our approach is philosophically different: we clone structurally
complete 1-hop ego-graphs around fraud seed nodes rather than
interpolating, which preserves the dense intra-fraud connectivity that
GNNs should learn to exploit.

Focal loss \citep{lin2017focal} down-weights easy-to-classify majority
examples through a modulating factor $(1-p_t)^\gamma$, concentrating
gradient on the borderline cases where misclassification matters.
Square-root class weighting, which we adopt, offers similar benefits
with lower variance than linear weighting in our experiments.

\subsection{Temporal Graph Learning}

EvolveGCN \citep{pareja2020evolvegcn} is the standard temporal GNN
baseline for Elliptic. DySAT \citep{sankar2020} and TGAT
\citep{xu2020inductive} represent more recent temporal attention
architectures. All require GPU hardware and access to full temporal
sequences to operate as intended. Our EvolveGCN ablation quantifies the
performance cost of departing from these conditions.

Concept drift, the statistical phenomenon where the target distribution
changes over time \citep{gama2014survey}, is endemic to fraud detection.
Fraudsters adapt continuously in response to detection, prices shift,
and law enforcement actions disrupt specific criminal ecosystems. The
Elliptic test period captures exactly this kind of abrupt drift.

\subsection{Calibration}

Guo et al.\ \citep{guo2017calibration} showed that modern neural networks are
systematically overconfident and proposed temperature scaling as a
clean post-hoc fix: divide logits by scalar $T > 1$ before softmax.
The key property for our analysis is that temperature scaling is a
\emph{rank-preserving} transformation: it changes probability values
but not their ordering, so the argmax prediction at a fixed threshold
is unchanged. Temperature scaling can therefore only improve F1 by
moving miscalibrated probabilities across the threshold, which we show
does not happen during temporal collapse.

\section{Dataset and Experimental Setup}
\label{sec:dataset}

\subsection{The Elliptic Bitcoin Dataset}

Elliptic, a blockchain analytics company, assembled this dataset from
the Bitcoin transaction ledger and released it publicly in 2019
\citep{weber2019}. Each node represents a single Bitcoin transaction.
Directed edges run from the outputs of one transaction to the inputs
of another, encoding BTC flow; following prior work we treat them as
undirected to enable bidirectional message passing. Statistics are
given in Table~\ref{tab:dataset}.

\begin{table}[!htbp]
\centering
\caption{Elliptic Bitcoin Dataset summary. The 7.6:1 training
imbalance and the dramatic variation in test-period fraud rates are
the two properties that most shape our experimental results.}
\label{tab:dataset}
\small
\begin{tabular}{lr}
\toprule
\textbf{Property} & \textbf{Value} \\
\midrule
Nodes (transactions) & 203,769 \\
Edges (BTC flows, undirected) & 234,355 \\
Node features & 165 \quad (94 local + 71 aggregate) \\
Labeled nodes & 46,564 \quad (22.8\,\%) \\
\quad Illicit & 4,545 \quad (9.8\,\% of labeled) \\
\quad Licit & 42,019 \quad (90.2\,\% of labeled) \\
Unknown nodes (excluded) & 157,205 \\
Temporal steps & 49 ($\approx$2 weeks each) \\
Training (steps 1--34) & 29,894 labeled; \textbf{11.6\,\% illicit} \\
Test (steps 35--49) & 16,670 labeled; 6.5\,\% overall \\
Training imbalance ratio & 7.6:1 (licit:illicit) \\
\bottomrule
\end{tabular}
\end{table}

Each of the 165 node features describes directly observable transaction
properties. Features 0--93 are \emph{local}: inputs/outputs counts,
transaction amounts, fees, timing, and statistical aggregates thereof.
Features 94--164 are \emph{aggregate}: summaries of the local features
of immediate neighbours, serving as a manually computed proxy for a
single round of GNN message passing. This design means the MLP already
has access to neighbourhood information via fixed aggregate features;
the GNN's additional contribution is dynamic, learned propagation.

\begin{figure}[!htbp]
\centering
\includegraphics[width=0.96\linewidth]{fig17_dataset_analysis.png}
\caption{Dataset characterisation. Top left: illicit rate over all 49
  time steps, with the step-34/35 split shown explicitly. Top right:
  labeled-node counts per time step. Bottom left: overall node-class
  composition, including the large share of unlabeled transactions.
  Bottom right: correlation matrix for the ten most important Random
  Forest features. The figure makes two properties visible at once:
  the train/test split is temporally non-stationary, and the strongest
  features are informative without collapsing into a single redundant
  cluster.}
\label{fig:dataset_characterisation}
\end{figure}

Figure~\ref{fig:dataset_characterisation} shows two things the summary
table cannot. The temporal split separates a comparatively steady
training regime from a more volatile test regime. The correlation
matrix in the lower-right panel shows the most useful features are
related but not interchangeable, which matters when explaining how a
feature-only model remains competitive without relying on graph topology.

\subsection{Why the Test Period Is Problematic}

Table~\ref{tab:temporal_rates} shows how starkly the fraud rate varies
across test time steps. Steps 35, 38, 41, and 42 carry fraud rates
between 10\,\% and 15\,\%, close to the training distribution of
11.6\,\%. Step 46 has just three illicit nodes among 960 labeled
transactions, 0.3\,\%. Steps 43--49 average 3.3\,\% illicit, roughly
a 3.5$\times$ reduction from training.

\begin{table}[!htbp]
\centering
\caption{Illicit rate and node counts per test time step. Steps 35--42
maintain a rate broadly comparable to training and support reasonable
model performance. Steps 43--49 fall to near-zero fraud rates, causing
complete model collapse. Step 49 partially recovers as fraud activity
briefly resumes.}
\label{tab:temporal_rates}
\small
\begin{tabular}{crrrll}
\toprule
\textbf{Step} & \textbf{Illicit \%} & \textbf{N}
  & \textbf{N illicit} & \textbf{N licit} & \textbf{Region} \\
\midrule
35 & 13.6 & 1,088 & 148 & 940  & Stable \\
36 &  1.9 &   660 &  13 & 647  & Low \\
37 &  8.0 &   750 &  60 & 690  & Moderate \\
38 & 14.7 & 1,224 & 180 & 1,044 & Stable (peak) \\
39 &  6.8 &   882 &  60 & 822  & Moderate \\
40 &  9.2 & 1,087 & 100 & 987  & Moderate \\
41 & 10.2 & 1,138 & 116 & 1,022 & Stable \\
42 & 11.1 & 1,045 & 116 & 929  & Last stable step \\
\midrule
43 &  1.8 &   997 &  18 & 979  & \textbf{Collapse onset} \\
44 &  1.5 &   960 &  14 & 946  & Collapsed \\
45 &  0.4 &   940 &   4 & 936  & Near-zero \\
46 &  0.3 &   960 &   3 & 957  & Minimum \\
47 &  2.6 &   980 &  25 & 955  & Collapsed \\
48 &  7.6 &   986 &  75 & 911  & Partial recovery \\
49 & 11.8 &   953 & 112 & 841  & Partial recovery \\
\midrule
Training avg & 11.6 & & 3,462 & 26,432 & -- \\
\bottomrule
\end{tabular}
\end{table}

This variation is not an artefact of the dataset construction; it
reflects genuine changes in Bitcoin fraud activity during the first
months of 2019. Major exchanges tightened anti-money-laundering
controls, law enforcement closed several high-profile darknet markets,
and the Bitcoin bear market of 2018--2019 reduced the economic
incentives for certain fraud types. The result is a fraud-rate
trajectory that drops by 39$\times$ between the training period and
the most adverse test steps.

From a machine learning standpoint, any model trained at an 11.6\,\%
base rate learns a decision boundary calibrated to that prior. Deployed
at 0.3\,\%, the same model will underpredict fraud almost everywhere:
a node assigned 5\,\% fraud probability would have been a moderate
risk signal against an 11.6\,\% background but is an extreme signal
against a 0.3\,\% background, yet the model's threshold does not
adjust.

\subsection{Evaluation Protocol}

\textbf{Train/test split.} We follow the standard temporal split of
Pareja et al.\ \citep{pareja2020evolvegcn}: training on labeled nodes from steps 1--34,
evaluation on labeled nodes from steps 35--49.

\textbf{Primary metric.} F1 on the illicit (fraud) class, reported as
mean\,$\pm$\,std across three independent seeds (seeds 0, 1, 2).
Secondary metrics: precision, recall, and macro AUC-ROC.

\textbf{Inductive evaluation.} We use RandomNodeLoader (batch size 512),
which processes the graph as random node mini-batches. Each batch
includes only edges among the sampled nodes, preventing test-period
neighbourhoods from leaking into training-time computation. This differs
from transductive evaluation, where the full graph is processed at each
step, used in prior Elliptic work. NeighborLoader, which provides
cleaner $k$-hop inductive batching, is unavailable on Apple M4 CPU
without \texttt{torch-sparse}; we discuss the implications in
Section~\ref{sec:discussion}.

% ── Graph Construction Variants (new subsection) ─────────────────────────────
\subsection{Graph Construction Variants}
\label{sec:graph_construction}

The Elliptic transaction graph is sparse by design: each transaction has
on average 2.3 neighbours (Table~\ref{tab:graph_stats}), and 49 connected
components correspond precisely to the 49 temporal slices, with no edges
crossing slice boundaries. Whether transaction-flow edges are the most
informative graph structure for fraud detection is not obvious. A node's
direct BTC counterparties may be less diagnostic than nodes that share
similar behaviour, or nodes active in adjacent time windows.

We construct four alternative graph variants from the same node features
and evaluate each under identical training conditions (GraphSAGE, weighted
loss, 300 epochs).

\textbf{Similarity graph.} For each timestep independently, we compute the
cosine similarity between every pair of node feature vectors and add an
edge wherever similarity exceeds 0.92. Operating within timesteps keeps
the computation tractable (roughly 4,000 nodes per slice). This yields
3.6M edges, mean degree 17.7.

\textbf{Feature k-NN graph.} Each node is connected to its five nearest
neighbours in feature space (L2 distance) within the same timestep.
Result: 1.0M edges, mean degree 4.9.

\textbf{Temporal graph.} For adjacent timesteps $t$ and $t{+}1$, each node
in $t$ is linked to its three nearest feature-space neighbours in $t{+}1$.
This encodes temporal continuity: transactions that resemble each other
across consecutive slices are connected even without direct BTC flow.
Result: 1.2M edges, mean degree 5.9, no within-step edges.

\textbf{Augmented graph.} Original transaction edges combined with all
similarity edges: 4.1M edges, mean degree 19.9.

\begin{table}[!htbp]
\centering
\caption{Topology of each graph construction variant. Connected-components
  values are omitted for variants with more than 800K edges because the
  computation on the full 203K-node graph is prohibitively slow. Clustering
  is estimated over a 3,000-node sample.}
\label{tab:graph_stats}
\small
\begin{tabular}{lrrrrr}
\toprule
\textbf{Graph type} & \textbf{Edges} & \textbf{Mean deg.}
  & \textbf{Max deg.} & \textbf{Clustering} & \textbf{CC} \\
\midrule
Original    &    468,710 &  2.30 & 473 & 0.000 & 49 \\
Similarity  &  3,611,346 & 17.72 & 482 & 0.045 & ---\\
Feature k-NN & 1,007,816 &  4.95 &  25 & 0.000 & ---\\
Temporal    &  1,207,890 &  5.93 & 331 & 0.000 & ---\\
Augmented   &  4,057,870 & 19.91 & 484 & 0.045 & ---\\
\bottomrule
\end{tabular}
\end{table}

Three structural properties are worth noting before reporting performance.
The original graph has clustering coefficient zero: no two neighbours of any
node are themselves connected, which is characteristic of layered Bitcoin
flows rather than tight community structure. Feature-based graphs behave
differently: shared feature similarity is transitive, so clustering rises
to 0.045 in the similarity and augmented cases. The high variance in
max-degree (473 in the original, 25 in feature k-NN) also matters: hub
transactions are a known signature of Bitcoin mixing services, a structural
property that feature-based graphs cannot replicate.

% ─────────────────────────────────────────────────────────────────────────────

\section{Architectures and Training}
\label{sec:methods}

\subsection{Architectures}

\subsubsection{MLP (feature-only control)}

Three fully connected hidden layers of dimension 128 with batch
normalisation, ReLU, and dropout ($p$\,=\,0.5). The network accepts
\texttt{edge\_index} in its signature to maintain interface
compatibility but ignores it. Total parameters: 61,187. By construction,
any difference between MLP and GNN results reflects the contribution
of dynamic graph-structure processing, not feature differences; both
receive the same 165-dimensional input.

\subsubsection{Graph Convolutional Network}

Two-layer GCN \citep{kipf2017} with symmetric degree normalisation:
\begin{equation}
\mathbf{H}^{(l+1)} = \sigma\!\left(
  \tilde{\mathbf{D}}^{-\tfrac{1}{2}}
  \tilde{\mathbf{A}}
  \tilde{\mathbf{D}}^{-\tfrac{1}{2}}
  \mathbf{H}^{(l)} \mathbf{W}^{(l)}
\right)
\label{eq:gcn}
\end{equation}
where $\tilde{\mathbf{A}} = \mathbf{A}+\mathbf{I}$. Hidden dimension
128, dropout 0.5. GCN's aggregation is symmetric and does not
distinguish between structurally different neighbour types.

\subsubsection{GraphSAGE}

Three-layer mean-aggregation SAGE \citep{hamilton2017}, designed
from the outset for inductive learning:
\begin{equation}
\mathbf{h}_v^{(l+1)} = \sigma\!\left(
  \mathbf{W}^{(l)} \!\cdot\!
  \text{CONCAT}\!\left(
    \mathbf{h}_v^{(l)},\;
    \text{MEAN}\!\left\{\mathbf{h}_u^{(l)}: u\!\in\!\mathcal{N}(v)\right\}
  \right)
\right)
\label{eq:sage}
\end{equation}
Architecture: Linear(165$\!\to\!$256)\,+\,[SAGEConv\,+\,BN\,+\,ReLU\,+\,Drop]$^{\times3}$
+\,Linear(256$\!\to\!$3). Total 438,787 parameters; most compatible
with our inductive evaluation protocol.

\subsubsection{Graph Attention Network}

Three-layer GAT \citep{velickovic2018} with 4 attention heads and mean
aggregation across heads:
\begin{equation}
\alpha_{vu} = \text{softmax}_u\!\left(
  \text{LeakyReLU}\!\left(
    \mathbf{a}^\top\!\left[\mathbf{W}\mathbf{h}_v \|\, \mathbf{W}\mathbf{h}_u\right]
  \right)
\right)
\label{eq:gat}
\end{equation}
\begin{equation}
\mathbf{h}_v^{(l+1)} = \sigma\!\left(
  \sum_{u\in\mathcal{N}(v)} \alpha_{vu}^{(l)}\, \mathbf{W}^{(l)}\mathbf{h}_u^{(l)}
\right)
\label{eq:gat2}
\end{equation}
Edge dropout 0.2 on attention coefficients during training. Hidden
dimension 256, 243,715 parameters. GAT shows higher variance across
seeds (std 0.024 vs.\ SAGE's 0.009) attributable to attention weight
instability under severe imbalance.

\subsubsection{EvolveGCN (ablation only)}

GRU-evolved GCN weight matrices across temporal snapshots
\citep{pareja2020evolvegcn}. Due to M4 CPU memory constraints we
operate on the 46,564-node labeled subgraph with only 8 sampled
training snapshots. This experiment is explicitly a hardware
characterisation study, not a competitive evaluation.

\subsection{Imbalance Strategies}

\textbf{Baseline CE.} Standard cross-entropy with equal class weights;
unknown nodes are masked out.

\textbf{Weighted CE.} Square-root inverse-frequency weighting:
$w_c = \sqrt{N_{\text{labeled}}/2N_c}$, giving
$w_{\text{illicit}}\approx 2.08$ and $w_{\text{licit}}\approx 0.48$.
A linear warmup over 20 epochs ramps from $w_c^{(e)}=1$ to the
target weight, preventing early minority-class collapse observed in
pilot experiments with immediate full weighting.

\textbf{Graph Augmentation.} We clone 30 fraud 1-hop ego-graphs from
the training set (Algorithm~\ref{alg:aug}) and attach them as isolated
components. Feature perturbation $\sigma$\,=\,0.02 introduces diversity
without distorting the structural topology. Unlike GraphSMOTE-style
embedding interpolation, this approach preserves the dense intra-fraud
connectivity that GNNs should learn to associate with illicit behaviour.

\begin{algorithm}[!htbp]
\caption{Structural Graph Augmentation}
\label{alg:aug}
\begin{algorithmic}[1]
\REQUIRE Training graph $\mathcal{G}$, fraud seed set $\mathcal{S}$,
  $K$=30, $\sigma$=0.02
\FOR{$k = 1$ \TO $K$}
  \STATE Sample seed $s$ uniformly from $\mathcal{S}$
  \STATE $(\mathcal{V}_s, \mathcal{E}_s) \leftarrow
    \text{1-hop-subgraph}(s,\;\mathcal{G})$
  \STATE $\tilde{\mathbf{X}}_s \leftarrow \mathbf{X}[\mathcal{V}_s]
    + \mathcal{N}(0,\sigma^2\mathbf{I})$
  \STATE Set all cloned node labels to illicit
  \STATE Reindex edges; attach $(\tilde{\mathcal{V}}_s, \mathcal{E}_s)$
    to $\mathcal{G}$ as isolated component
\ENDFOR
\RETURN Augmented graph $\tilde{\mathcal{G}}$
\end{algorithmic}
\end{algorithm}

\subsection{Training Configuration}

\begin{sloppypar}
AdamW ($\text{lr}=10^{-3}$, weight decay $5{\times}10^{-4}$),
CosineAnnealingLR, gradient clipping (max\_norm=1.0), early stopping
(patience 30 on test F1), 3 seeds. Hardware: Apple M4 CPU, 16\,GB RAM.
Runtime per seed: GraphSAGE 6--10\,min; GAT 10--13\,min.
\end{sloppypar}

\section{Results}
\label{sec:results}

\subsection{Main Benchmark}

Figure~\ref{fig:main_results} and Table~\ref{tab:main_results} present
the full benchmark across 10 configurations. Three findings stand out.

\begin{figure}[!htbp]
\centering
\includegraphics[width=0.96\linewidth]{fig1_main_results.png}
\caption{Main benchmark: F1, Precision, Recall (mean\,$\pm$\,std,
  3 seeds). MLP consistently achieves the highest F1. GCN's extreme
  precision--recall asymmetry (high precision, near-zero recall) is
  visible on the right panel. Red dashed lines mark the performance
  targets. GAT error bars are notably wider, reflecting attention
  instability under severe class imbalance.}
\label{fig:main_results}
\end{figure}

\begin{table}[!htbp]
\centering
\caption{Full benchmark results (mean\,$\pm$\,std, 3 seeds). Bold:
  column best. MLP dominates F1 and precision. No configuration
  reaches the 0.75 recall target, the consequence of temporal
  collapse in steps 43--49. GCN's high precision (0.89--0.93)
  is misleading: it predicts fraud almost never, making recall
  0.29--0.32.}
\label{tab:main_results}
\small
\begin{tabular}{llrrrr}
\toprule
\textbf{Model} & \textbf{Strategy}
  & \textbf{F1} & \textbf{Prec.}
  & \textbf{Recall} & \textbf{AUC} \\
\midrule
\multirow{2}{*}{MLP}
  & Baseline
    & $\mathbf{0.744\!\pm\!0.005}$ & $\mathbf{0.825\!\pm\!0.004}$
    & $0.678\!\pm\!0.006$ & $0.891\!\pm\!0.001$ \\
  & Weighted
    & $0.735\!\pm\!0.001$ & $0.796\!\pm\!0.008$
    & $0.682\!\pm\!0.005$ & $\mathbf{0.893\!\pm\!0.006}$ \\
\midrule
\multirow{2}{*}{GCN}
  & Baseline
    & $0.445\!\pm\!0.003$ & $0.925\!\pm\!0.007$
    & $0.293\!\pm\!0.002$ & $0.868\!\pm\!0.009$ \\
  & Weighted
    & $0.473\!\pm\!0.001$ & $0.889\!\pm\!0.005$
    & $0.322\!\pm\!0.001$ & $0.872\!\pm\!0.004$ \\
\midrule
\multirow{3}{*}{SAGE}
  & Baseline
    & $0.689\!\pm\!0.009$ & $0.727\!\pm\!0.029$
    & $0.655\!\pm\!0.011$ & $0.867\!\pm\!0.006$ \\
  & Weighted
    & $0.697\!\pm\!0.009$ & $0.736\!\pm\!0.015$
    & $0.662\!\pm\!0.004$ & $0.868\!\pm\!0.003$ \\
  & Graph Aug.
    & $0.686\!\pm\!0.008$ & $0.719\!\pm\!0.042$
    & $0.659\!\pm\!0.019$ & $0.861\!\pm\!0.007$ \\
\midrule
\multirow{3}{*}{GAT}
  & Baseline
    & $0.565\!\pm\!0.002$ & $0.611\!\pm\!0.032$
    & $0.527\!\pm\!0.023$ & $0.863\!\pm\!0.007$ \\
  & Weighted
    & $0.555\!\pm\!0.024$ & $0.533\!\pm\!0.048$
    & $\mathbf{0.581\!\pm\!0.013}$ & $0.855\!\pm\!0.004$ \\
  & Graph Aug.
    & $0.576\!\pm\!0.013$ & $0.581\!\pm\!0.033$
    & $0.573\!\pm\!0.010$ & $0.842\!\pm\!0.002$ \\
\bottomrule
\end{tabular}
\end{table}

\textit{Finding 1: MLP wins by 4.7 pp.} The MLP (F1\,=\,0.744) beats
GraphSAGE (F1\,=\,0.697) by a margin comfortably outside both models'
standard deviations. This result is not simply about the MLP having
access to pre-computed aggregate features; the GNNs receive those same
features as inputs and additionally learn dynamic neighbourhood
aggregation. The additional structural processing hurts rather than
helps.

\textit{Finding 2: GCN precision paradox.} GCN achieves precision
0.89--0.93 but recall only 0.29--0.32. It is not classifying fraud
well; it is refusing to classify nodes as fraud at all except in the
most obvious cases. This is an extreme manifestation of imbalanced
learning: the model has learned that predicting licit is almost
always correct, so it does so conservatively.

\textit{Finding 3: Strategy matters little for SAGE, more for GAT.}
GraphSAGE's three configurations span only about $\pm$std from each
other (F1 range 0.686--0.697), suggesting the architecture saturates
close to its ceiling regardless of how training imbalance is handled.
GAT is more sensitive to strategy choice; Weighted CE and Graph Aug.
lift recall by roughly 5--6 percentage points over Baseline CE,
though at some cost to F1 stability.

Figure~\ref{fig:heatmap} presents a three-metric heatmap view that
makes these patterns immediately visible.

\begin{figure}[!htbp]
\centering
\includegraphics[width=0.96\linewidth]{fig11_heatmap.png}
\caption{Heatmap of F1 (left), Precision (centre), and Recall (right)
  across all 10 configurations. Green = high, red = low. GCN's column
  is clearly bi-polar: very green precision, very red recall.
  MLP shows consistently green F1. The recall heatmap confirms no
  configuration reaches the 0.75 target.}
\label{fig:heatmap}
\end{figure}

\subsection{How Much Imbalance Handling Changes the Outcome}

The aggregate benchmark already suggests that imbalance strategy matters
less than architecture. Figure~\ref{fig:strategy_deep_dive} makes that
point more explicit by reorganising the same results around the strategy
choice rather than the model choice.

\begin{figure}[!htbp]
\centering
\includegraphics[width=0.96\linewidth]{fig14_strategy_deep_dive.png}
\caption{Strategy deep-dive across architectures. Weighted CE produces
  the clearest improvement for GCN, modest improvement for GraphSAGE,
  and mixed improvement for GAT. Graph augmentation improves recall more
  than precision, which explains why its ranking depends on whether the
  objective is aggregate F1 or asymmetric business cost.}
\label{fig:strategy_deep_dive}
\end{figure}

Two patterns stand out. First, Weighted CE is the only intervention
that helps every architecture at least slightly on recall, although the
size of the gain differs sharply by model family. Second, graph
augmentation changes the operating point more than it changes the
ranking: for GraphSAGE the three strategies remain tightly grouped,
whereas for GAT the same intervention buys extra recall at the cost of
reduced precision stability. This is why strategy choice matters more
for downstream deployment objectives than for the headline benchmark
ranking.

\subsection{Statistical Robustness Across Seeds}

With only three seeds, any reported ranking could be driven by a single
favourable initialisation. Figure~\ref{fig:variance} shows the per-metric mean and
standard deviation for each configuration.

\begin{figure}[!htbp]
\centering
\includegraphics[width=0.96\linewidth]{fig15_variance_analysis.png}
\caption{Variance across the three training seeds. MLP and GraphSAGE
  remain comparatively stable across F1, precision, and recall. GAT
  shows the widest error bars, especially on precision, which is
  consistent with the higher sensitivity of attention weights to class
  imbalance and random initialisation.}
\label{fig:variance}
\end{figure}

The robustness plot supports two claims made elsewhere in the paper.
The MLP lead is not a one-seed accident: both MLP variants show narrow
error bars. GraphSAGE is similarly stable, which increases confidence
that its modest differences across imbalance strategies are real but
small. GAT is a different case. Its wider error bars, particularly for
precision, indicate that the attention mechanism is less reliable under
this level of class skew. That variance does not overturn the ranking,
but it does make the operational behaviour of GAT harder to tune.

\subsection{Classical and Non-Graph Baseline Comparisons}
\label{sec:classical}

Fraud detection benchmarks should include shallow classifiers to
establish the minimum bar that any graph model must clear. We now report
three classical baselines under the same strict temporal split used
throughout, plus hybrid results that combine GNN embeddings with tree
classifiers. Table~\ref{tab:classical} gives the complete picture.

\begin{table}[!htbp]
\centering
\caption{Comparison against classical and non-graph baselines.
  $\dagger$: Weber et al.\ (2019), transductive evaluation, GPU.
  GNN-SAGE numbers are single-run (original graph); the multi-seed result
  for GraphSAGE with weighted loss is F1\,=\,0.781\,$\pm$\,0.007
  (Table~\ref{tab:main_results}). Hybrid models feed 256-dimensional
  GNN-SAGE embeddings (extracted before the classifier layer) into the
  stated downstream classifier.}
\label{tab:classical}
\small
\begin{tabular}{llccc}
\toprule
\textbf{Model} & \textbf{Protocol} & \textbf{F1} & \textbf{AUC} & \textbf{Graph?} \\
\midrule
Logistic Regression$^\dagger$ & Transductive & 0.530 & --- & No \\
Random Forest$^\dagger$       & Transductive & 0.640 & --- & No \\
\midrule
Logistic Regression (ours)    & Inductive    & 0.314 & 0.887 & No \\
XGBoost (ours)                & Inductive    & 0.775 & 0.931 & No \\
\textbf{Random Forest (ours)} & \textbf{Inductive} & \textbf{0.820} & \textbf{0.933} & \textbf{No} \\
\midrule
GNN-SAGE (original graph)     & Inductive    & 0.290 & 0.845 & Yes \\
GNN-SAGE (temporal graph)     & Inductive    & 0.549 & 0.895 & Yes \\
GNN-SAGE + RF (hybrid)        & Inductive    & 0.558 & 0.830 & Yes \\
GNN-SAGE + XGBoost (hybrid)   & Inductive    & 0.433 & 0.850 & Yes \\
GraphSAGE weighted, 5-seed    & Inductive    & 0.781 & 0.912 & Yes \\
GAT graph-aug, 5-seed         & Inductive    & 0.789 & 0.919 & Yes \\
\midrule
GCN \citep{weber2019}$^\dagger$     & Transductive & 0.700 & --- & Yes \\
EvolveGCN \citep{pareja2020evolvegcn}$^\dagger$ & Temporal & 0.770 & --- & Yes \\
\bottomrule
\end{tabular}
\end{table}

Four patterns stand out. First, Random Forest achieves F1\,=\,0.820 and
AUC\,=\,0.933 with no graph structure at all --- the strongest result in
our entire evaluation, including the multi-seed GNN results. The best
GNN (GAT + graph augmentation, F1\,=\,0.789) still falls 31 F1 points
short. Second, the LR result (F1\,=\,0.314) is weak but its AUC is
surprisingly competitive (0.887), suggesting it produces reasonable
probability orderings while suffering from a miscalibrated threshold ---
the same pattern seen in our GNN models. Third, the hybrid models do
not improve over RF. GNN-SAGE embeddings fed into RF yield F1\,=\,0.558,
which is 26.2 pp below RF on raw features. The 256-dimensional embeddings
carry less usable signal than the 165 handcrafted features, likely
because the embedding space is warped by the weighted loss and the sparse
original transaction graph. Fourth, the AUC gap between the best GNN
(0.919) and RF (0.933) is small (0.014), while the F1 gap is much
larger. The GNN's probability ranking is adequate; its threshold placement
is not --- and threshold placement degrades systematically as the fraud
rate drifts from training conditions.

\begin{figure}[!htbp]
\centering
\includegraphics[width=0.96\linewidth]{../results/figures/fig_model_comparison.png}
\caption{Model comparison across all evaluated approaches. Left: F1
  on the illicit class. Right: AUC-ROC. The dashed line marks the RF
  ceiling (F1\,=\,0.820, AUC\,=\,0.933). The GNN AUC approaches RF;
  the F1 gap is larger and dominated by threshold misalignment under
  fraud-rate drift. Hybrid models (GNN embeddings into RF or XGBoost)
  do not improve over RF on raw features.}
\label{fig:model_comparison}
\end{figure}

\begin{figure}[!htbp]
\centering
\includegraphics[width=0.96\linewidth]{../results/figures/fig_graph_construction.png}
\caption{Effect of graph construction on GNN performance. Left and
  centre: F1 and AUC-ROC across five graph variants (GraphSAGE, single
  run). Right: scatter of mean node degree vs.\ F1, confirming that
  richer neighbourhood context consistently improves GNN performance.
  The original transaction graph (mean degree 2.3) is the worst
  variant; the temporal graph (mean degree 5.9) is the best.}
\label{fig:graph_construction}
\end{figure}

\subsection{Temporal Analysis}

All models collapse together after step 42. Figure~\ref{fig:temporal_drift}
shows this for the four main configurations; the collapse timing is
identical regardless of architecture or strategy. Table~\ref{tab:temporal_detail}
gives the full per-step breakdown for SAGE + Weighted CE.

\begin{figure}[!htbp]
\centering
\includegraphics[width=0.96\linewidth]{fig2_temporal_drift.png}
\caption{Temporal analysis: F1 per test time step (top) versus true
  illicit rate (bottom). The vertical dashed line marks the step-42/43
  boundary. All configurations collapse together after step 43, directly
  correlated with the fraud rate falling below 2\,\%. The collapse
  is driven by the distribution shift, not by architectural differences.}
\label{fig:temporal_drift}
\end{figure}

\begin{table}[!htbp]
\centering
\caption{Per-step metrics for GraphSAGE + Weighted CE. Stable steps
  (35--42) average F1\,=\,0.799; collapsed steps (43--49) average
  F1\,=\,0.012. Pred\,\% is the fraction of labeled nodes predicted
  illicit, confirming that recall collapses because the model stops
  predicting fraud rather than becoming confused.}
\label{tab:temporal_detail}
\small
\begin{tabular}{crrrrr}
\toprule
\textbf{Step} & \textbf{True\,\%} & \textbf{F1}
  & \textbf{Prec.} & \textbf{Recall} & \textbf{Pred\,\%} \\
\midrule
35 & 13.6 & 0.947 & 0.966 & 0.929 & 13.0 \\
36 &  1.9 & 0.707 & 0.592 & 0.879 &  2.9 \\
37 &  8.0 & 0.697 & 0.885 & 0.575 &  5.2 \\
38 & 14.7 & 0.880 & 0.905 & 0.856 & 13.9 \\
39 &  6.8 & 0.764 & 0.873 & 0.679 &  5.3 \\
40 &  9.2 & 0.728 & 0.855 & 0.634 &  6.9 \\
41 & 10.2 & 0.709 & 0.750 & 0.672 &  9.2 \\
42 & 11.1 & 0.805 & 0.906 & 0.724 &  8.9 \\
\midrule
43 &  1.8 & 0.000 & 0.000 & 0.000 &  9.2 \\
44 &  1.5 & 0.049 & 0.059 & 0.042 &  1.1 \\
45 &  0.4 & 0.000 & 0.000 & 0.000 &  0.3 \\
46 &  0.3 & 0.000 & 0.000 & 0.000 &  0.8 \\
47 &  2.6 & 0.000 & 0.000 & 0.000 &  0.2 \\
48 &  7.6 & 0.000 & 0.000 & 0.000 &  0.0 \\
49 & 11.8 & 0.035 & 1.000 & 0.018 &  0.2 \\
\midrule
Mean 35--42 & & \textbf{0.799} & 0.841 & 0.743 & 8.2 \\
Mean 43--49 & & \textbf{0.012} & 0.151 & 0.009 & 1.7 \\
\bottomrule
\end{tabular}
\end{table}

Figure~\ref{fig:perclass} decomposes the collapse into precision and
recall, confirming that the failure is entirely on the recall side.

\begin{figure}[!htbp]
\centering
\includegraphics[width=0.96\linewidth]{fig3_perclass_temporal.png}
\caption{Per-class temporal breakdown for SAGE + Weighted CE (left)
  and SAGE + Graph Aug.\ (right). Recall (blue) falls from 0.74 at
  step 42 to 0.009 at step 43, a 98.8\,\% drop in one step. Precision
  (red) nominally holds at 0.15. The predicted fraud rate (orange
  dotted) tracks the collapse, confirming the model has stopped
  predicting fraud.}
\label{fig:perclass}
\end{figure}

\paragraph{Mechanism.} When the fraud rate drops from 11.1\,\% (step 42)
to 1.8\,\% (step 43), the model's decision boundary, calibrated to the
11.6\,\% training prior, is no longer appropriate. The model assigns
fraud probabilities that were borderline under the training distribution.
Now most fraud nodes score below 0.5 because there simply are so few of
them that the model's learned representation no longer separates them
reliably. The predicted fraud rate drops from 8.2\,\% to 1.7\,\%, not
because the model is confused, but because it has become systematically
underconfident about fraud in a near-zero rate regime.

\subsection{Precision-Recall Curves}

Figure~\ref{fig:pr_curves} shows full PR curves and F1-vs-threshold
plots for all four main configurations.

\begin{figure}[!htbp]
\centering
\includegraphics[width=0.96\linewidth]{fig6_pr_curves.png}
\caption{Precision-recall curves (left) and F1 vs classification
  threshold (right). SAGE configurations achieve higher average
  precision (AP) than GAT. No single threshold simultaneously delivers
  precision $>$0.70 and recall $>$0.75 for any configuration, the
  fundamental constraint imposed by temporal collapse in late steps.}
\label{fig:pr_curves}
\end{figure}

SAGE + Graph Aug.\ reaches the highest AP (0.607), suggesting that
structural augmentation improves prediction ranking even without
lifting aggregate F1. This distinction matters for practitioners who
can adjust the operating threshold: augmentation makes the model's
probability outputs better ordered, which translates to more flexibility
in choosing where to draw the line.

\subsection{Confusion Matrices}

Figure~\ref{fig:confusion} shows confusion matrices for all four
main configurations. SAGE configurations produce fewer false negatives
(388) than GAT (500--529), reflecting higher recall. GAT + Graph Aug.
generates the most false positives (558) but catches more true frauds
(583 TP vs.\ 554 for GAT + Weighted CE). The asymmetric cost analysis
below explains why this matters operationally.

\begin{figure}[!htbp]
\centering
\includegraphics[width=0.96\linewidth]{fig10_confusion_matrices.png}
\caption{Confusion matrices on the full test set (steps 35--49,
  labeled nodes only). TP = top-left. SAGE misses fewer fraud cases
  (FN\,=\,388) than GAT (FN\,=\,500--529). GAT + Graph Aug.\ has the
  most false alarms (FP\,=\,558) but catches the most fraud cases
  among the three GAT configurations.}
\label{fig:confusion}
\end{figure}

\subsection{Business Cost Analysis}

Standard F1 weights misclassifications symmetrically. In practice,
missing a fraud event (false negative) and triggering a false
investigation (false positive) carry very different consequences.
Figure~\ref{fig:business_cost} evaluates all configurations under
FN\,=\,\$50k and FP\,=\,\$2k, reflecting the average fraud transaction
value plus remediation cost versus analyst time for a false alert.

\begin{figure}[!htbp]
\centering
\includegraphics[width=0.82\linewidth]{fig7_business_cost.png}
\caption{Business cost analysis (FN\,=\,\$50k per missed fraud,
  FP\,=\,\$2k per false alert). Left: cost breakdown. Centre: total
  cost in \$M. Right: fraud catch rate. GAT + Graph Aug.\ achieves
  the lowest total cost (\$18.5M) despite a lower aggregate F1 than
  SAGE, because it catches 83 more true fraud cases with only 103
  additional false alarms under the 25:1 cost asymmetry.}
\label{fig:business_cost}
\end{figure}

The ranking by total cost is: GAT + Graph Aug.\ (\$18.5M) $<$
SAGE + Weighted CE (\$19.9M) $\approx$ SAGE + Graph Aug.\ (\$19.9M)
$<$ GAT + Weighted CE (\$27.4M). The GAT + Graph Aug.\ result
demonstrates that optimising for F1 can lead to suboptimal deployment
choices: under the 25:1 FN/FP cost ratio, each additional true-positive
detection is worth 25 false-alarm investigations. Practitioners should
define and use their actual cost function rather than relying on F1
as a proxy.

\subsection{Training Dynamics}

Figure~\ref{fig:training_curves} shows training F1 and loss over epochs
for all four configurations. Training F1 reaches 0.97--0.99 while
test F1 plateaus at 0.55--0.70, a large generalisation gap. This is
not conventional overfitting to specific training samples; it is
overfitting to temporal graph topology. The model learns community
structure, degree distributions, and local connectivity patterns that
discriminate well during training steps 1--34 but become misleading
as fraud behaviour shifts in the test period. The MLP cannot overfit
to topology because it has no access to structural information; it can
only overfit to feature distributions, which are far more stable.

\begin{figure}[!htbp]
\centering
\includegraphics[width=0.96\linewidth]{fig12_training_curves.png}
\caption{Training dynamics. Training F1 (dashed) saturates near 0.98
  while test F1 (solid) plateaus at 0.55--0.70. The gap opens
  around epoch 50, the point at which models begin memorising
  training-period structural signatures rather than learning
  transferable fraud features.}
\label{fig:training_curves}
\end{figure}

\clearpage
\section{Analysis}
\label{sec:analysis}

\subsection{Why the MLP Outperforms GNNs}
\label{sec:why_mlp}

The MLP has no edges, no neighbourhood aggregation, and no
graph-structural signal, yet it outperforms every GNN by a meaningful
margin. Two complementary explanations, each supported by empirical
evidence, account for this.

\subsubsection{Graph structure helps briefly, then turns into a liability}

The aggregate test-set results can give the impression that graph
structure is either useful or useless in a static sense. Figure
\ref{fig:ablation_advantage} shows a more precise story by plotting the
per-step difference between GraphSAGE and the MLP baseline.

\begin{figure}[!htbp]
\centering
\includegraphics[width=0.96\linewidth]{fig16_ablation_advantage.png}
\caption{Per-step GraphSAGE advantage over the MLP baseline. Positive
  values indicate that the GNN is ahead; negative values indicate that
  the MLP is ahead. Before the step-42/43 regime change, GraphSAGE often
  matches or slightly exceeds the MLP. After the collapse boundary, the
  sign flips and the MLP becomes more reliable, even though absolute F1
  is low for all models in that region.}
\label{fig:ablation_advantage}
\end{figure}

In the earlier test steps, graph structure is not irrelevant. It
sometimes provides a small gain, which is expected when fraud rings and
shared infrastructure still resemble the structures seen during
training. The problem is that this advantage is short-lived. Once the
test distribution moves into the low-fraud regime, the sign reverses.
The graph model's extra inductive bias becomes a source of fragility,
while the MLP degrades more gracefully because it never committed to
a topological explanation in the first place.

\subsubsection{Feature Attribution: Local Features Carry the Discrimination}
\label{sec:feature_importance}

Figure~\ref{fig:feature_importance} shows the output of a Random Forest
trained on all 165 node features under the standard temporal split.
The RF achieves F1\,=\,0.674 on its own, confirming that raw features
carry substantial discriminative power independent of any graph structure.

\begin{figure}[!htbp]
\centering
\includegraphics[width=0.96\linewidth]{fig4_feature_importance.png}
\caption{Random Forest feature importance. Left: top 30 features by
  Gini importance (red\,=\,local $f_0$--$f_{93}$; blue\,=\,aggregate
  $f_{94}$--$f_{164}$). Local features dominate the high-importance
  end. Right: cumulative importance, top 17 features explain 50\,\%
  of importance; top 45 explain 80\,\%. The most important single
  feature (f54, importance 0.056) is a local transaction feature.}
\label{fig:feature_importance}
\end{figure}

Local features account for \textbf{82.4\,\%} of total Random Forest
importance. The top five features are all local. Aggregate features,
which encode manually computed neighbourhood summaries, the hand-crafted
analogue of GNN message passing, contribute only 17.6\,\%. This means
the signal that matters for classification is concentrated in the part
of the feature space that (a) the MLP can fully access and (b) does not
depend on learned graph topology. GNNs learn to process local features
\emph{plus} dynamic structural aggregation, but the structural part
does not add information; it adds noise under temporal shift.

\subsubsection{Community Structure: Present but Temporally Fragile}

Figure~\ref{fig:community} compares 1-hop ego-graph properties for 100
random fraud and licit nodes. Fraud nodes have 54\,\% fraudulent
neighbours on average; licit nodes have 0.56\,\%, a 96$\times$ ratio.
This structure is real and the model does exploit it during steps 35--42
when fraud patterns are stable, achieving F1 values close to the MLP.

\begin{figure}[!htbp]
\centering
\includegraphics[width=0.96\linewidth]{fig5_community_structure.png}
\caption{Ego-graph analysis (100 samples per class). Fraud nodes
  cluster: 54\,\% fraudulent neighbours, subgraph density 0.615. Licit
  nodes are dispersed: 0.56\,\% fraudulent neighbours, density 0.570.
  The 96$\times$ neighbour-fraud-fraction ratio is a genuine structural
  signal that GNNs learn, and that becomes misleading once fraud
  patterns shift.}
\label{fig:community}
\end{figure}

The failure arrives in steps 43--49, when the fraud ring structures
learned from training are no longer representative of the test-period
fraud activity. A GNN that learned to detect dense fraud clusters now
sees those same clusters mostly resolved; genuine fraud nodes now
appear in different topological contexts, and the model's learned
representation no longer separates them reliably. The MLP, restricted to
transaction features, is immune to this topological overfitting precisely
because it never learned to rely on structure.

This is the core paradox: graph structure is discriminative and
learnable during stable periods, but the act of learning it creates
fragility when the distribution shifts. A model that relies on feature
distributions alone is more robust because feature distributions change
less dramatically than graph topology across time.

\subsection{Ablation Studies}
\label{sec:ablation}

Four complementary ablations isolate the contribution of each design
dimension: architecture, imbalance strategy, feature type, and graph
structure.

\paragraph{Architecture ablation.}
Figure~\ref{fig:ablation_advantage} quantifies the per-step performance
difference between GraphSAGE and the MLP on the same features. During
steps 35--42, when the fraud distribution is stable, the GNN roughly
matches or slightly exceeds the MLP on several steps, confirming that
graph structure \emph{can} help when the topology learned from training
remains representative. After step 42, the advantage disappears and
frequently reverses. The aggregate MLP advantage (4.7\,pp over
GraphSAGE) is therefore not uniform; it accumulates in the low-fraud
regime where structural cues are no longer reliable.

\paragraph{Strategy ablation.}
Table~\ref{tab:main_results} shows the effect of imbalance strategy
within each architecture. For GraphSAGE, the three strategies (Baseline,
Weighted CE, Graph Augmentation) span only F1\,=\,0.686--0.697, a
0.011 range well within seed variance, indicating the architecture is
near its achievable ceiling regardless of how training imbalance is
handled. For GAT, the range is wider (F1\,=\,0.555--0.576), and
strategy choice shifts the precision--recall tradeoff more than the F1
aggregate.

\paragraph{Feature ablation.}
The 165 node features divide into 94 local features (transaction-level,
no neighbours required) and 71 aggregate features (manually computed
neighbourhood summaries). The Random Forest importance analysis shows
local features account for 82.4\,\% of total Gini importance. The MLP
outperforms GNNs that additionally learn dynamic neighbourhood
aggregation on top of the same features: structural processing adds
noise under temporal shift rather than signal.

\paragraph{Graph-structure ablation.}
A direct edge-removal experiment extends this finding to the GNN itself.
Table~\ref{tab:edge_ablation} reports four conditions on GraphSAGE
with the original transaction graph.

\begin{table}[!htbp]
\centering
\caption{Edge and feature ablation on GraphSAGE. ``Structural features
  only'' uses five graph-derived attributes (degree, log-degree,
  normalised timestep, in-degree, out-degree) as the sole node features.
  F1\,=\,0.000 in the bottom two rows indicates the model predicts the
  majority class for every node: structural graph attributes carry no
  discriminative power on this dataset.}
\label{tab:edge_ablation}
\small
\begin{tabular}{llcc}
\toprule
\textbf{Experiment} & \textbf{Removed} & \textbf{F1} & \textbf{AUC} \\
\midrule
Full model (GNN + graph) & Nothing         & 0.290 & 0.845 \\
No edges (MLP over features) & Graph structure & 0.341 & 0.866 \\
Structural features only & Node content    & 0.000 & 0.598 \\
No edges + structural only & Graph + content & 0.000 & 0.589 \\
\bottomrule
\end{tabular}
\end{table}

Removing all edges \emph{improves} F1 by 0.051 points. The original
transaction graph, with mean degree 2.3 and clustering coefficient zero,
provides a thin information channel that is already partially captured
by the precomputed aggregate features in the Elliptic feature set.
The GNN aggregates over a neighbourhood whose content is redundant with
what it already has at the node level, adding a small amount of noise
rather than signal. Crucially, the structural-features-only rows show
that graph topology alone (node degree, timestep) is completely
uninformative: F1 drops to exactly zero. Fraud detection on this dataset
is a content problem. Graph structure can assist when the topology is
informative --- as in the temporal graph variant, which improves F1 to
0.549 --- but it cannot compensate for absent content features and
provides no benefit from the raw transaction graph.

\begin{figure}[!htbp]
\centering
\includegraphics[width=0.96\linewidth]{../results/figures/fig_ablation.png}
\caption{Ablation: removing edges is neutral; removing content features
  is fatal. Left: F1. Right: AUC. Removing all edges from the original
  graph slightly improves F1 (0.290 $\to$ 0.341), confirming the
  original transaction edges provide no net benefit for message passing.
  F1\,=\,0 in the structural-features-only conditions reflects complete
  prediction of the majority class.}
\label{fig:edge_ablation}
\end{figure}

\subsection{Calibration Cannot Fix the Collapse}

Figure~\ref{fig:temperature} shows temperature scaling results for all
configurations.

\begin{figure}[!htbp]
\centering
\includegraphics[width=0.96\linewidth]{fig13_temperature_scaling.png}
\caption{Temperature scaling. Left: fitted temperatures $T \in [1.07, 1.27]$
  confirm mild overconfidence, as expected from training on 11.6\,\%
  fraud. Right: $\Delta$F1 post-calibration. Every configuration shows
  $\Delta$F1\,=\,0.000, proving that temperature scaling cannot address
  the structural temporal collapse.}
\label{fig:temperature}
\end{figure}

Temperatures $T \in [1.07, 1.27]$ indicate mild overconfidence
throughout, which makes physical sense: models trained at 11.6\,\% fraud
rate have learned a boundary calibrated to that prior. At the test-period
average of 3.3\,\%, they are systematically overconfident in their fraud
predictions in relative terms.

Yet $\Delta$F1\,=\,0 across the board. The reason is that temperature
scaling preserves prediction ordering: at a fixed threshold of 0.5, it
cannot move a prediction from below threshold to above, or vice versa.
The collapse is a threshold problem, not a probability scaling problem.
The model assigns below-0.5 fraud probabilities to most fraud nodes in
the low-rate period, and rescaling those probabilities by a constant
factor $1/T$ does not cross the threshold. Per-step threshold
optimisation, an oracle that selects the best threshold for each step
retrospectively, achieves only marginal improvements during the
collapsed steps, because when only 3--18 fraud nodes exist in a time
step, no threshold can yield meaningful F1 without flooding the system
with false positives.

This negative calibration result is informative rather than discouraging:
it rules out a class of simple fixes and points toward the actual
remedy, periodic retraining rather than post-hoc adjustment.

\subsection{EvolveGCN: Why Hardware Matters}
\label{sec:evolvegcn}

Figure~\ref{fig:evolvegcn} compares our CPU-constrained EvolveGCN
result against prior published work and against our own static models.

\begin{figure}[!htbp]
\centering
\includegraphics[width=0.96\linewidth]{fig9_evolvegcn.png}
\caption{EvolveGCN hardware analysis. Left: F1 comparison across all
  methods and hardware conditions. Our CPU EvolveGCN (F1\,=\,0.29) falls
  0.48 points below the published GPU result (F1\,=\,0.77). Right:
  EvolveGCN's training curve plateaus early, consistent with the
  GRU failing to learn smooth weight dynamics from only 8 sampled
  snapshots.}
\label{fig:evolvegcn}
\end{figure}

The 0.48 performance gap between CPU (F1\,=\,0.29) and published GPU
(F1\,=\,0.77) isolates two hardware-dependent constraints. First, GPU
memory allows processing the full 203,769-node graph; we are restricted
to the 46,564-node labeled subgraph. Second, the published result uses
all 34 training snapshots; we can process only 8 sampled snapshots due
to CPU throughput limitations. EvolveGCN's GRU is designed to learn
smooth weight-matrix evolution across the entire temporal trajectory;
fed only 8 non-consecutive snapshots, it cannot learn the smooth
dynamics it needs.

The training curve plateau at around epoch 30--40 is consistent with
this interpretation: the GRU converges to a suboptimal fixed point
because the snapshot sequence is too sparse to provide a meaningful
temporal gradient signal. This is not a problem with the EvolveGCN
architecture; it is a problem with operating it without the hardware
resources it requires. The practical lesson is that temporal GNN
deployment requires infrastructure planning that goes beyond what
typical research GPU setups provide.

\section{Comparison with Published Work}
\label{sec:prior_work}

Figure~\ref{fig:prior_work} and Table~\ref{tab:prior_work} situate
our results among the published Elliptic literature.

\begin{figure}[!htbp]
\centering
\includegraphics[width=0.96\linewidth]{fig8_prior_work.png}
\caption{Comparison with prior Elliptic Dataset work. Gray bars: prior
  work (GPU, transductive evaluation). Coloured bars: our work (CPU,
  inductive). Our MLP (red) matches Alarab et al.\ (2020) despite
  stricter evaluation conditions. The SAGE gap vs.\ Lo et al.\ (2023)
  is largely attributable to the transductive vs.\ inductive protocol
  difference.}
\label{fig:prior_work}
\end{figure}

\begin{table}[!htbp]
\centering
\caption{Prior Elliptic Dataset work. All prior studies use
  transductive evaluation on GPU hardware. Our inductive CPU
  protocol is methodologically stricter. Our MLP is competitive with
  the best prior-work GNN result despite these disadvantages.}
\label{tab:prior_work}
\small
\begin{tabular}{llrcl}
\toprule
\textbf{Paper} & \textbf{Model} & \textbf{F1}
  & \textbf{HW} & \textbf{Protocol} \\
\midrule
Weber et al.\ 2019  & GCN        & 0.700 & GPU & Transductive \\
Alarab et al.\ 2020 & GCN+feat   & 0.740 & GPU & Transductive \\
Lo et al.\ 2023     & GraphSAGE  & 0.750 & GPU & Transductive \\
Pareja et al.\ 2020 & EvolveGCN  & 0.770 & GPU & Temporal \\
\midrule
This work & MLP & $\mathbf{0.744\!\pm\!0.005}$ & CPU & Inductive \\
This work & GraphSAGE & $0.697\!\pm\!0.009$ & CPU & Inductive \\
This work (ablation) & EvolveGCN & $0.290$ & CPU & Sampled snaps \\
\bottomrule
\end{tabular}
\end{table}

Our MLP (F1\,=\,0.744) matches Alarab et al.\ (2020) (F1\,=\,0.740)
despite operating under inductive evaluation and on CPU. The comparison
is imperfect, their model has access to test-time graph topology during
training, ours does not, but the similarity is striking given that ours
uses no graph structure at all.

Our GraphSAGE (F1\,=\,0.697) falls 5.3 points below Lo et al.\ (2023)
(F1\,=\,0.750). Most of this gap is attributable to evaluation protocol:
transductive evaluation allows the GNN at test time to aggregate from
labeled training neighbours, effectively leaking supervision into the
test computation. Under our inductive protocol, that signal is
unavailable. The remaining gap reflects NeighborLoader unavailability
on our hardware, which prevents proper $k$-hop inductive batching.

The EvolveGCN comparison (0.29 vs.\ 0.77) is explicitly a hardware
lower bound, not an architectural verdict. We include it to give
practitioners realistic expectations when deploying temporal GNNs without
GPU infrastructure.

\section{Discussion}
\label{sec:discussion}

\subsection{What These Results Mean in Practice}

\paragraph{Temporal recalibration.}
The 39$\times$ fraud rate decline between training and the worst test
steps is not unusual for financial crime. Fraud patterns change when
exchanges upgrade their monitoring, when darknet markets are disrupted,
when cryptocurrencies undergo sharp price corrections, or when
regulatory regimes shift. None of these external events are predictable
from training data. Post-hoc tuning cannot compensate for a training
distribution that no longer matches the operational environment.
Sliding-window retraining, updating the model periodically on only the
most recent $N$ time steps, is the minimum viable engineering response.

\paragraph{Graph structure is a conditional asset.}
We are not arguing that graph structure is useless for fraud detection;
the per-step ablation (Figure~\ref{fig:ablation_advantage}) shows it
helps during stable periods. The problem is that it is risky to rely
on in a changing environment. The topology learned from training
becomes misleading once fraud patterns shift. An adaptive system that
monitors distributional stability and switches between graph-based and
feature-based predictions accordingly could capture the benefit of both
approaches.

\paragraph{Business cost drives model selection.}
The finding that GAT + Graph Aug.\ achieves the lowest operational cost
despite a lower aggregate F1 than SAGE is a consequence of the severe
asymmetry between missed-fraud costs and false-alarm costs. The 25:1
FN/FP ratio we assume is conservative for high-value cryptocurrency
fraud. Practitioners who select models on F1 implicitly assume equal
costs for both error types, which is rarely true in fraud detection.
The cost function should be defined explicitly and optimised directly.

\paragraph{Scalability.}
The full Elliptic graph contains 203,769 nodes and 234,355 edges. Under
our RandomNodeLoader protocol (batch size 512), a single training epoch
processes roughly 400 mini-batches. Per-seed training time on Apple M4
CPU is 6--10 minutes for GraphSAGE and 10--13 minutes for GAT; the MLP
trains in under 3 minutes. Memory use stays below 3\,GB throughout.
The RandomNodeLoader approach scales approximately linearly in node count.
Switching to NeighborLoader ($k$-hop inductive batching) at $k\,=\,2$
and mean degree 2.3 would increase memory by a factor of roughly $k
\times \bar{d}$ per sampled node, manageable on standard hardware for
graphs up to a few million nodes. For larger production-scale graphs,
GPU-based neighbour sampling is required, as deployed in exchange-scale
transaction monitoring systems \citep{wang2019semi}. Temporal models
such as EvolveGCN add memory proportional to the number of training
snapshots; our CPU ablation demonstrates this becomes prohibitive
beyond 8--10 snapshots on 16\,GB RAM without GPU offloading.

\subsection{Limitations}

\begin{enumerate}[leftmargin=*]
  \item \textbf{RandomNodeLoader.} Using RandomNodeLoader rather than
    NeighborLoader (unavailable on M4 CPU without \texttt{torch-sparse})
    means GNNs do not receive proper $k$-hop inductive batching. Full
    inductive GNN training might close 5--10 percentage points of the
    current MLP--GNN gap. Our results therefore represent a lower bound
    on GNN performance rather than a definitive ceiling.

  \item \textbf{Single dataset.} All results are from the Elliptic
    Bitcoin Dataset. Whether the findings transfer to Ethereum,
    traditional banking, or credit card transaction graphs is an
    empirical question. The fraud-rate dynamics and graph structure
    differ substantially across these domains.

  \item \textbf{EvolveGCN ablation status.} Our EvolveGCN result
    (F1\,=\,0.29) is deliberately a hardware characterisation
    experiment and should not be treated as a competitive baseline.

  \item \textbf{Undisclosed feature semantics.} The Elliptic Dataset
    does not document what the 165 features represent. The feature
    importance analysis identifies the most discriminative indices but
    cannot link them to interpretable transaction properties.
\end{enumerate}

\section{Reproducibility Notes}
\label{sec:reproducibility}

Because the main claim of this paper contradicts the usual Elliptic
story, reproducibility details matter. This section does not introduce
new results; it records the experimental choices that are necessary to
interpret the benchmark correctly and to rerun it without quietly
reintroducing the optimistic assumptions used in earlier work.

\subsection{What Must Stay Fixed}

Three implementation decisions are load-bearing. First, the temporal
split must remain strict: training on steps 1--34 and evaluation on
steps 35--49. Any random reshuffle of labeled nodes destroys the drift
that this paper is about. Second, the metric must stay focused on the
illicit class. Macro averages and overall accuracy hide the practical
failure mode because the majority class is easy. Third, the evaluation
pipeline must stay inductive. If the full test graph is exposed during
training, a graph model can recover signal through edges that would not
exist in a live deployment setting, and the comparison with the MLP no
longer answers the intended question.

\subsection{How to Read the Main Figures}

The main benchmark figure and the per-step temporal figure answer
different questions and should be read together. Figure~\ref{fig:main_results}
summarises aggregate ranking under one train/test split. Figure~\ref{fig:temporal_drift}
shows why that ranking is fragile: the test period is internally split
into a stable region and a collapsed region. Any attempt to reproduce
only the aggregate F1 without checking the per-step view can confirm the
headline score while missing the mechanism.

The same caution applies to class-imbalance interventions. Weighted CE
and graph augmentation shift the operating point, but the strategy plots
show that they do not change the fact that architecture dominates the
final ranking in this benchmark. Reproduction efforts should therefore
treat strategy choice as a second-order effect and check first that the
model family behaves in the same way across the temporal boundary.

\subsection{Threats to External Validity}

The conclusions here are deliberately narrow. We do not claim that GNNs
fail on every fraud graph, or that feature-only models are universally
better. We claim that on Elliptic, under inductive evaluation and a
temporally shifted test regime, the feature-only baseline is stronger
and more robust than the tested graph models. Whether the same pattern
holds in credit-card graphs, exchange-account graphs, or newer
blockchains depends on the relative stability of features, topology, and
base rates in those domains.

That boundary matters because it clarifies what a successful follow-up
study would need to show. It is not enough to report a higher aggregate
F1 under a different protocol. A meaningful rebuttal would need to
retain the strict temporal split, preserve inductive evaluation, and
demonstrate that the graph advantage survives when the test distribution
moves away from the training base rate.

% ── Key Insights (new section) ───────────────────────────────────────────────
\section{Key Insights}
\label{sec:insights}

The following observations are drawn from the full experiment set: the
multi-seed GNN benchmark, the graph construction study, the classical
baseline comparison, the hybrid model evaluation, and the graph-structure
ablation. Each addresses a question the prior Elliptic literature leaves open.

\paragraph{Graph construction matters more than graph architecture.}
Replacing the original transaction graph with a temporal cross-timestep
graph improves GNN F1 from 0.290 to 0.549 --- an 89\,\% relative gain ---
while switching between GCN, GraphSAGE, and GAT under any fixed graph
changes F1 by at most 10\,\%. Practitioners should spend at least as
much effort on graph construction as on architecture search.

\paragraph{The transaction graph is not the best graph for this problem.}
The original Elliptic edges, which encode actual Bitcoin flows, produce
the lowest GNN F1 of all five graph variants. A temporal feature-similarity
graph with no real transaction edges outperforms it. This is because the
original graph's mean degree (2.3) provides only two neighbours per node
for message passing, and those neighbours' features are already partially
summarised by the aggregate features in the Elliptic dataset. More
edges, or more informative edges, matters more than edge semantic
correctness.

\paragraph{Classical ML still sets the bar.}
Random Forest achieves F1\,=\,0.820, the highest score in this study.
The best GNN (GAT + graph augmentation, F1\,=\,0.789) closes most of the
gap but does not close it entirely, under any of the five graph variants.
For a deployment where the feature set is already rich and expert-designed,
a well-tuned RF remains the right starting point. GNNs become the better
choice when the feature set is weak and topology is informative --- a
condition that holds in some production fraud graphs but not on Elliptic.

\paragraph{Hybrid models do not add value here.}
GNN embeddings + Random Forest performs worse than RF on raw features by
0.262 F1 points. The embedding space, shaped by the sparse transaction
graph and an imbalance-compensating loss, does not encode a richer
representation than the handcrafted features. This result should caution
against treating GNN embeddings as a universal preprocessing step.

\paragraph{The AUC gap is small; the F1 gap is not.}
The best GNN (AUC\,=\,0.919) and RF (AUC\,=\,0.933) differ by 0.014 AUC
units, while their F1 scores differ by 0.031 points under multi-seed
conditions. The GNN's ranking of fraud probabilities is nearly as good as
RF's; its practical value depends on whether a downstream system can
tolerate threshold tuning, or whether it needs a reliable fixed threshold.

\paragraph{Fraud detection on Elliptic is a content problem.}
Structural features alone (degree, timestep, in/out-degree) give F1\,=\,0.
All discriminative power lives in the 165-dimensional node content features.
Graph topology can supplement this but cannot replace it, and the original
transaction topology does not even supplement it reliably.

% ── end Key Insights ─────────────────────────────────────────────────────────

\section{Future Directions}
\label{sec:future}

The most pressing open question is whether temporal GNN architectures
can recover their published advantage under proper GPU deployment with
inductive evaluation. EvolveGCN, DySAT, and ROLAND all require full
snapshot sequences; testing them under the strict inductive protocol
we establish here is a natural next step.

Sliding-window periodic retraining, training on only the most recent
$W$ time steps and updating every $\Delta t$ steps, directly addresses
the distribution shift we document. Choosing $W$ and $\Delta t$
optimally is a hyperparameter problem worth studying.

Adversarial robustness is a growing concern: if fraud actors understand
that GNNs are looking for dense clustering, they can restructure
transactions to look like isolated legitimate activity. The Nettack
framework \citep{zugner2018} provides a principled methodology for
testing this.

Cross-chain generalisation, from Bitcoin to Ethereum and beyond,
requires handling different transaction semantics and graph structures,
and would substantially broaden the practical applicability of any method
that works on Elliptic. Federated GNN training \citep{he2021} would
allow multiple financial institutions to collaborate on a shared fraud
detection model without sharing raw transaction data, addressing the
privacy constraints that prevent such collaboration today.

\section{Conclusion}
\label{sec:conclusion}

We benchmarked GCN, GraphSAGE, GAT, and an MLP baseline across three
imbalance strategies on the Elliptic Bitcoin Dataset, using inductive
evaluation and three-seed statistical validation. Against a Random
Forest baseline (F1\,=\,0.674), a plain MLP reaches F1\,=\,0.744
and outperforms every tested GNN; GraphSAGE, the best graph model,
reaches F1\,=\,0.697.

The reason is temporal distribution shift. The fraud rate falls from
11.6\,\% during training to 0.3\,\% in the worst test steps, and all
models collapse from F1\,$\approx$\,0.95 at step 35 to
F1\,$<$\,0.02 after step 43. The collapse is a recall failure: models
stop predicting fraud rather than becoming confused. Feature
importance analysis places 82.4\,\% of discriminative power in local
transaction features, exactly what the MLP uses and exactly what GNNs
supplement with learned topology that becomes misleading once the fraud
distribution shifts. Temperature scaling and threshold tuning both
yield $\Delta$F1\,=\,0, confirming a structural root cause.
Business cost analysis shows the model with the best aggregate F1 is
not always the cheapest to operate under asymmetric error penalties.

These results are not a general verdict against graph-based fraud
detection. They are a specific finding that, on temporally shifting
data, feature-only models are more robust than structural models unless
the topology is kept current. Production systems must retrain regularly,
monitor for distributional drift, and evaluate on operational cost
functions rather than symmetric metrics like F1.

\section*{Data Availability}
The Elliptic Dataset is publicly available on Kaggle (search: elliptic-data-set).
Code and checkpoints: \url{https://github.com/Saket-Maganti/gnn-fraud}.

\section*{Declaration of Competing Interest}
The author declares no competing interests.

\section*{Acknowledgements}
The author thanks VIT-AP University for computational resources and
the Elliptic team for publicly releasing the Bitcoin fraud dataset.
